\subsection{Explicit polynomial witness clock}
The source bookkeeping supplies, at source step $j$, a total primary-register
value at most $2+j$, excess header gap at most $j$, first spare value 1, and
remaining spare-list width at most $75+3 j$. The gap estimate uses the separately
checked fact that every nonhalting source successor has address-overflow flag 0;
without that fact the displayed estimate would not follow merely from the
decrement semantics. Initially there are 26 unit spare blocks. An increment
appends at most one such block, while a positive decrement contributes at most
one to the excess gap.

Write a source-outcome duration vector as
$(c_0,c_1,c_2,\ldots,c_{10},c_{11},c_{12})$ and put
$m=\max_{2\leq i\leq 10}c_i$. All coefficients are nonnegative, so the associated
physical duration is at most
\[
 \bigl(c_0+2 m+c_{11}+75 c_{12}\bigr)
 +\bigl(c_1+m+3 c_{12}\bigr)j.
\]
The finite coefficient audit gives maxima 694 and 18 for these two expressions.
Consequently the cumulative MP clock through $n$ source steps satisfies
\[
 b(n)\leq 58066+\sum_{j<n}(694+18 j)
       =58066+685 n+9 n^2=: \Theta(n).
\]
The parent starts after 57922 physical steps and satisfies the same coefficient
envelope, hence its clock is bounded by the same $\Theta$. Source halt is
extended by a dummy absorbing segment of duration 1; every ordinary segment has
positive duration. Both clocks are therefore at least the source-step index.

With the separately established exact startup and uniform source-frame
simulation, define the two clocks by primitive recursion along the common
extended source trajectory. If either physical machine has halted by time $t$,
then it is halted at its clock value for source index $t$. The corresponding
source state is therefore the absorbing halt state, and the other machine is
halted by its own clock value, at most $\Theta(t)$. These primitive recursive
witness translations supply the relative arithmetic halting equivalence.
They do not discharge the separate proof-calculus calibration to
$\operatorname{Con}(\mathrm{PA})$.
